%! Independent Events and Unions — Unit 4, Lesson 4.6 — Student Packet Cover Sheet
\documentclass[10pt]{article}
\usepackage{apstats-article}
\usepackage{apstats-boxes}
\usepackage{ltablex}
\keepXColumns

\begin{document}

% Full-bleed navy banner
\begin{tikzpicture}[remember picture, overlay]
  \fill[navy] (current page.north west)
    rectangle ([yshift=-0.9in]current page.north east);
\end{tikzpicture}
\vspace{-0.8in}

\begin{center}
  {\color{white}\textbf{\LARGE AP Statistics}}\\[4pt]
  {\color{skymid}\large\textbf{Unit 4: Probability, Random Variables, and Probability Distributions}}\\[6pt]
  {\color{white}\textbf{\large Lesson 4.6 \quad Independent Events and Unions of Events}}
\end{center}

\vspace{0.22in}
\namedateperiod
\vspace{0.18in}

\begin{learningtargetbox}
\small
\begin{enumerate}[leftmargin=1.5em, itemsep=5pt]
  \item \ldots{} determine whether two events are independent by checking whether
        $P(A|B) = P(A)$ or equivalently $P(A \cap B) = P(A) \cdot P(B)$.
  \item \ldots{} calculate $P(A \cup B)$ using the addition rule
        $P(A \cup B) = P(A) + P(B) - P(A \cap B)$.
  \item \ldots{} explain the difference between independent events and mutually
        exclusive events, and why mutually exclusive events with positive probabilities
        cannot be independent.
\end{enumerate}
\end{learningtargetbox}

\vspace{0.14in}

\begin{tocbox}
\small
\renewcommand{\arraystretch}{1.5}
\begin{tabularx}{\linewidth}{@{} c l X r @{}}
\rowcolor{sky}
\textbf{\#} & \textbf{Component} & \textbf{Description} & \textbf{Score} \\
\midrule
1 & Warm-Up        & Spiral review: conditional probability and mutually exclusive    & \blank{1.2cm} \\
2 & Guided Notes   & Independence definition; addition rule; ME vs.\ independent      & \blank{1.2cm} \\
3 & Group Activity & School program $\times$ satisfaction survey table                & \blank{1.2cm} \\
4 & Exit Ticket    & School preferences two-way table                                 & \blank{1.2cm} \\
5 & Homework       & Streaming service $\times$ age group table                       & \blank{1.2cm} \\
\midrule
  & \multicolumn{2}{r}{\textbf{Total}} & \blank{1.2cm} \\
\end{tabularx}
\end{tocbox}

\vspace{0.14in}

\begin{remindbox}
\small
The \textbf{two big ideas of Lesson 4.6}: independence and the addition rule.

\vspace{6pt}
\begin{tabularx}{\linewidth}{@{} >{\centering\arraybackslash\columncolor{goldbg}}p{0.06\linewidth}
                                    X @{}}
\textbf{\large!} &
\textbf{Independence test: compare $P(A|B)$ to $P(A)$.}\enspace
If they are equal, knowing $B$ did not change the probability of $A$ --- the
events are independent. If they differ, the events are dependent. \\[6pt]
\textbf{\large!} &
\textbf{Addition rule: subtract the intersection to avoid double-counting.}\enspace
$P(A \cup B) = P(A) + P(B) - P(A \cap B)$. For mutually exclusive events,
$P(A \cap B) = 0$, so the rule simplifies. \\[6pt]
\textbf{\large!} &
\textbf{Mutually exclusive $\neq$ independent --- they are almost opposites.}\enspace
If two events are mutually exclusive and both have positive probability, they
are \emph{dependent}: knowing one occurred guarantees the other did not. \\
\end{tabularx}
\end{remindbox}

\end{document}
